
\documentclass[11pt]{article}
\usepackage[letterpaper,margin=0.75in]{geometry}
\usepackage[T1]{fontenc}
\usepackage[utf8]{inputenc}
\usepackage{lmodern}
\usepackage{array}
\usepackage{longtable}
\usepackage{enumitem}
\usepackage{xcolor}
\usepackage{amssymb}
\usepackage{xurl}
\usepackage[hidelinks]{hyperref}
\setlength{\parindent}{0pt}
\setlength{\parskip}{6pt}
\emergencystretch=3em
\setlist[itemize]{leftmargin=1.2em,itemsep=2pt,topsep=2pt}
\setlist[enumerate]{leftmargin=1.4em,itemsep=2pt,topsep=2pt}
\renewcommand{\arraystretch}{1.25}
\newcommand{\checkbox}{$\Box$\hspace{0.4em}}
\newcommand{\ok}{\textbf{Success criteria:} }
\newcommand{\sourcebasis}{\section*{Official Source Basis}\begin{itemize}
\item Authentik Docker Compose installation: \url{https://docs.goauthentik.io/install-config/install/docker-compose/}
\item Authentik reverse proxy guidance: \url{https://docs.goauthentik.io/install-config/reverse-proxy/}
\item Authentik configuration reference: \url{https://docs.goauthentik.io/install-config/configuration/}
\item Authentik first steps: \url{https://docs.goauthentik.io/install-config/first-steps/}
\item Authentik backup and restore: \url{https://docs.goauthentik.io/sys-mgmt/ops/backup-restore/}
\item Authentik upgrade guidance: \url{https://docs.goauthentik.io/install-config/upgrade/}
\item Authentik monitoring: \url{https://docs.goauthentik.io/sys-mgmt/ops/monitoring/}
\item Authentik example flows: \url{https://docs.goauthentik.io/add-secure-apps/flows-stages/flow/examples/flows/}
\item Authentik service accounts: \url{https://docs.goauthentik.io/sys-mgmt/service-accounts/}
\item Authentik certificates: \url{https://docs.goauthentik.io/sys-mgmt/certificates/}
\end{itemize}}

\title{Authentik Docker Compose Deployment Guide}\author{Prepared from official Authentik documentation}\date{May 6, 2026}\begin{document}\maketitle

\section*{Purpose}
This guide provides a task-oriented Docker Compose deployment procedure for Authentik based only on the official Authentik documentation. It assumes a single host, Docker Compose v2, a reverse proxy for HTTPS, and persistent PostgreSQL-backed configuration.

\section*{Prerequisites}
\begin{itemize}
\item Host with at least 2 CPU cores and 2 GB of RAM.
\item Docker Compose v2 or Podman with Compose support.
\item DNS name for Authentik if it will be accessed through a reverse proxy.
\item Reverse proxy capable of HTTPS and WebSocket upgrade forwarding.
\item Administrative control over the host filesystem, Compose project, and backup process.
\end{itemize}

\section*{Procedure 1 - Create the Installation Directory}
\begin{verbatim}
sudo mkdir -p /opt/authentik
sudo chown "$USER":"$USER" /opt/authentik
cd /opt/authentik
\end{verbatim}
\ok The operator is in a dedicated Authentik installation directory.

\section*{Procedure 2 - Download the Official Compose File}
Linux:
\begin{verbatim}
wget https://docs.goauthentik.io/compose.yml
\end{verbatim}
macOS or systems preferring curl:
\begin{verbatim}
curl -O https://docs.goauthentik.io/compose.yml
\end{verbatim}
\ok The directory contains the official \texttt{compose.yml} file.

\section*{Procedure 3 - Generate Required Secrets}
For a fresh installation, generate the PostgreSQL password and Authentik secret key into \texttt{.env}:
\begin{verbatim}
echo "PG_PASS=$(openssl rand -base64 36 | tr -d '\n')" >> .env
echo "AUTHENTIK_SECRET_KEY=$(openssl rand -base64 60 | tr -d '\n')" >> .env
\end{verbatim}
\ok The \texttt{.env} file contains \texttt{PG\_PASS} and \texttt{AUTHENTIK\_SECRET\_KEY}.

\section*{Procedure 4 - Optional Configuration}
Enable Authentik error reporting only if desired:
\begin{verbatim}
echo "AUTHENTIK_ERROR_REPORTING__ENABLED=true" >> .env
\end{verbatim}
Expose HTTP and HTTPS on alternative host ports if needed:
\begin{verbatim}
cat >> .env <<'EOF'
COMPOSE_PORT_HTTP=80
COMPOSE_PORT_HTTPS=443
EOF
\end{verbatim}
Apply later configuration changes with:
\begin{verbatim}
docker compose up -d
\end{verbatim}
Verify the applied configuration when needed:
\begin{verbatim}
docker compose run --rm worker ak dump_config
\end{verbatim}
\ok Optional settings are documented in \texttt{.env} and can be verified through \texttt{ak dump\_config}.

\section*{Procedure 5 - Start Authentik}
\begin{verbatim}
docker compose pull
docker compose up -d
\end{verbatim}
\ok The Authentik server, worker, and PostgreSQL database are started by Docker Compose.

\section*{Procedure 6 - Complete Initial Setup}
Open the initial setup URL. Include the trailing slash:
\begin{verbatim}
http://<your server's IP or hostname>:9000/if/flow/initial-setup/
\end{verbatim}
Set the default \texttt{akadmin} password.
\ok The default administrator can log in to Authentik.

\section*{Procedure 7 - Configure the Reverse Proxy}
The reverse proxy must pass these upstream headers:
\begin{itemize}
\item \texttt{X-Forwarded-Proto}
\item \texttt{X-Forwarded-For}
\item \texttt{Host}
\item \texttt{Connection: Upgrade} and \texttt{Upgrade: WebSocket} for HTTP/1.1 WebSocket endpoints
\end{itemize}
When the proxy is not accessing Authentik from a trusted private range, configure trusted proxy CIDRs:
\begin{verbatim}
AUTHENTIK_LISTEN__TRUSTED_PROXY_CIDRS=<proxy-cidr-list>
\end{verbatim}
\ok Authentik receives correct protocol, host, client IP, and WebSocket forwarding context.

\section*{Procedure 8 - Add the First Application and User}
\begin{enumerate}
\item Log in as an administrator and open the Admin interface.
\item Navigate to \texttt{Applications > Applications}.
\item Create an application and provider pair.
\item Record provider values required by the target application, such as client ID, client secret, slug, redirect URI, and signing key.
\item Navigate to \texttt{Directory > Users}.
\item Create the first non-admin user.
\item Add bindings or policies when access should be restricted to selected users or groups.
\end{enumerate}
\ok The first protected application and user are available.

\section*{Procedure 9 - Back Up the Deployment}
Back up PostgreSQL using native PostgreSQL tooling such as \texttt{pg\_dump}, \texttt{pg\_dumpall}, or continuous archiving. Exclude system databases \texttt{template0} and \texttt{template1}. Back up static directories when they are used locally:
\begin{itemize}
\item \texttt{/data}
\item \texttt{/certs}
\item \texttt{/custom-templates}
\item \texttt{/blueprints}
\end{itemize}
\ok The database and deployment-specific static content can be restored.

\section*{Procedure 10 - Upgrade Authentik}
Before upgrading, read the release notes, back up PostgreSQL, and do not skip required major/minor upgrade sequence steps. Then run:
\begin{verbatim}
wget -O compose.yml https://docs.goauthentik.io/compose.yml
docker compose pull
docker compose up -d
\end{verbatim}
Verify the version from \texttt{Dashboards > Overview}. Ensure Authentik outposts are upgraded with the core instance.
\ok The deployment is upgraded and the running version is verified.

\sourcebasis

\end{document}
